\documentclass[12pt]{article}
\usepackage{verbatim}
\usepackage[dvips]{epsfig}
\usepackage{color}
\usepackage{url}
\usepackage[colorlinks=true]{hyperref}

\begin{document}

\section*{GENESIS: Documentation}

{\bf Related Documentation:}
% start: userdocs-tag-replace-items related-do-nothing
% end: userdocs-tag-replace-items related-do-nothing


Introduction
GENESIS 3.0 (G-3) conforms to the CBI federated software architecture [1]. The modular nature of the system is designed to facilitate the interface of different databases that significantly enhance the user and developer experience. These databases fall into two broad categories. One is organized around an ideal user workflow [2] that lies at the core of simulator use. The other supports a multi-level documentation system that informs user and developer concerns. Here we introduce and explain how this system supports scientific communication using community models.
Methods
The GENESIS documentation system is supported by a database that comprises seven levels: Level 1: Introductory material, Level 2: User guides and documentation, Level 3: Automated regression testing and use cases, Level 4: Technical guide specification, Level 5: Algorithm documentation, Level 6: API documentation, and Level 7: Inline source code documentation


use the existing documentation system from a viewing angle of computational model publications.

Level 1: Descriptions of evolution of computational neuroscience and
general neuroscience.

Level 2: Reviews about community models in computational neuroscience.

Level 3: Narratives of computational model publications with an
intended audience of non-technical readers.

Level 4: Narratives of the biological mechanisms explained by
computational models for people with as main interest computational
modeling.

Level 5: Implementation and algorithm documents.

Level 6: API documentation.

Level 7: Source code and implementation.


\section*{New Document}

The three above headings are included in all documentation by default.

\begin{itemize}
\item[]  {\bf The Documentation Identifier:} This may be replaced by a descriptor of the type of documentation you have developed, for example {\bf Introduction} or {\bf Tutorial}, etc.
\item[]{\bf Related Documentation Heading:} Followed by a hyperlink list of any related documentation.
\item[] {\bf New Document Heading:} Should be replaced with the title of your document. {\bf Note:} In your document (the \LaTeX\,\,source documentation file--{\it NewDocument.tex}) all text below the {\bf New Document} heading should be replaced with your own content. See \href{../document-create/document-create.tex}{\bf Create\,Document} for further information.
\end{itemize}

\subsection*{Help with \LaTeX\,Documentation}

\subsubsection*{Text color}

Text in your document can be colorized  with the command
\begin{verbatim}
    \textcolor{red}{text}
\end{verbatim}
for example \textcolor{red}{\bf This is IMPORTANT}, \textcolor{green}{green} and \textcolor{blue}{blue} are also recognized arguments.

\subsubsection*{Hyperlinks}

Two classes of hyperlink exist in the GENESIS Documentation System:

\begin{itemize}
\item[{\bf A.}]{\bf Static Hyperlinks:} 
   \begin{itemize}
      \item{\bf Local links:} These are produced with the command:
\begin{verbatim}
   \href{../NewDocument/NewDocument.tex}{NewDocument}}
\end{verbatim}
for example \href{../NewDocument/NewDocument.tex}{NewDocument} references this document.
Note that these links may not work in your \LaTeX\,{\tt pdf} file viewer, but should be testable with the \href{http://get.adobe.com/reader/}{Acrobat Reader} or \href{http://www.adobe.com/products/acrobatpro/tryout.html}{Acrobat Professional} packages.
      \item {\bf Remote links:} These use the {\tt href} command in a slightly different way:
\begin{verbatim}
    \href{http://www.genesis-sim.org/documentation}{GENESIS Documentation}
\end{verbatim}
for example see \href{http://www.genesis-sim.org/documentation}{GENESIS Documentation}. Note that these links do not wrap automatically at the end of a line even outside of the {\tt verbatim} environment.
\end{itemize}

\item[{\bf B.}]{\bf Reference Hyperlinks:} See \href{../document-create/document-create.tex}{Document Creation} for further details on managing reference hyperlinks.

\end{itemize}

\subsubsection*{Figures}

Figures can easily be added with the following code:

\begin{verbatim}
\begin{figure}[h]
   \centering
   \includegraphics[scale=0.6]{figures/dummyfig.eps}
   \caption{{\bf A Dummy Figure:} Example of Latex code to incorporate
      a figure into documentation.}
   \label{fig:df-1}
\end{figure}
\end{verbatim}

\begin{figure}[h]
  \centering
   \includegraphics[scale=0.5]{figures/dummyfig.eps}
\caption{{\bf A Dummy Figure:} Example of \LaTeX\,\,\,code to incorporate a figure into documentation.}
  \label{fig:df-2}
\end{figure}

\bibliographystyle{plain}
\bibliography{../tex/bib/g3-refs.bib}

\end{document}
